\section{Admissibility and meta-theory}\label{sec:admissibility-meta-theory}

% Mathematical job:
%   - Introduce honest bookkeeping as an admissibility regime.
%   - Replace flat generator independence with typed non-collapse.
%   - State the level-profile architecture (behavioral / verification /
%     structural-categorical) as common meta-theory.
%   - State the forgetting / selected-lift architecture (preserved invariants,
%     loss records, non-uniqueness, lower-to-higher non-determinacy).
%   - State the activation-threshold architecture (below-threshold collapse,
%     threshold crossing, no-arbitrary-enrichment).
%   - State instrument-relative visibility and suppression.
%   - State empirical-bridge admissibility (threshold-dependent, level-indexed,
%     instrument-visible; not empirical realization).
%
% Governing sources:
%   workspace/meta/lens_lab/papers/paper2/theorem_sequence_claims.md
%   workspace/meta/lens_lab/papers/paper2/meta_theory_synthesis/paper2_meta_theory_synthesis.md
%   workspace/meta/lens_lab/papers/paper2/meta_theory_synthesis/admissibility_nonclaim_closure.md
%   workspace/meta/lens_lab/papers/paper2/scope_and_nonclaims.md
%   docs/papers/Tsiokos_2026_Six_Birds_No_Go_Theorems_for_Audited_Emergence.tex
%     (NG_LADDER_IDEM, NG_LADDER_BOUNDED_INTERFACE, NG_PROTOCOL_TRAP,
%     NG_MACRO_CLOSURE_DEFICIT)
%   paper/prior_literature_context.md (Tier A1, A2, A3)
%
% Overclaim risks:
%   - Collapsing typed non-collapse into global independence.
%   - Turning admissibility remarks into philosophy.
%   - Claiming instrument-family completeness (this is future work only).
%   - Claiming lossless forgetting or unique lifts without proof.

The admissibility regime is honest bookkeeping, elevated from a
per-description condition on \(D \in \FATCD{}\) to a meta-theoretic
discipline governing which claims about admissible descriptions enter
the theorem chain. Under honest bookkeeping, the flat
primitive-independence question is replaced by typed non-collapse,
level transitions are controlled by forgetting and selected lifts,
thresholds govern activation, instruments govern visibility, and
empirical-bridge claims are admitted only as threshold-dependent,
level-indexed, instrument-visible assertions. Seven theorem schemas
formalize the regime and preserve the global nonclaims of
Remark~\ref{rem:atlas-global-nonclaims}.

\subsection{Honest bookkeeping as an admissibility regime}\label{subsec:honest-bookkeeping}

The audit condition of Definition~\ref{def:audited} records enough
data for each claim, promotion, translation, lift, forgetting,
visibility, and empirical-bridge assertion to be interrogated on its
own terms. The following proposition states that these record
requirements jointly constitute the admissibility regime, in a spirit
analogous to the axiomatic specification discipline of Hoare
logic~\citep{Hoare1969}.

\begin{proposition}[Honest bookkeeping is an admissibility regime]\label{prop:honest-bookkeeping}
Honest bookkeeping is a necessary part of the admissibility regime for
claims about a description \(D \in \FATCD{}\). Concretely, any claim that
enters the theorem chain is admissible only if every assertion below
carries an explicit record in \(\Audit{D}\):
\begin{itemize}[topsep=2pt,itemsep=1pt]
  \item the claim itself, with its statement, hypotheses, and scope;
  \item every promotion, with its source level, target level, added data,
    losses, and preserved nonclaims;
  \item every translation, with its source context, target context,
    preserved data, suppressed data, and claim strength;
  \item every selected lift, marked selected and non-unique;
  \item every forgetting map, with its record of losses;
  \item every visibility claim, with both its visible and its suppressed
    data;
  \item every empirical-bridge assertion, with its substrate, observable,
    instrument, level map, threshold-witness relation, and
    suppressed-structure nonclaims.
\end{itemize}
Accordingly, unrecorded claims, implicit promotions, silent translations,
asserted-but-unselected lifts, lossless forgetting without proof,
visibility overreads, and bare empirical-realization assertions are all
inadmissible.
\end{proposition}

\begin{remark}[Why the regime is necessary]\label{rem:why-honest-bookkeeping}
The bookkeeping requirement is a structural check, not a stylistic
one. Without explicit records, passive observation can be misread as
\Psix{}, mere packaging as \Pone{} descent, and route failure as
\Pthree{} mismatch. The regime prevents those misreadings from
entering the theorem chain.
\end{remark}


\subsection{Typed non-collapse}\label{subsec:typed-non-collapse}

The regime does not include flat primitive independence --- the claim
that \Pone{}, \ldots, \Psix{} are independent generators of a global
algebra. Its meta-theoretic counterpart here is typed non-collapse.

\begin{proposition}[Typed non-collapse]\label{prop:typed-non-collapse}
Under honest bookkeeping (Proposition~\ref{prop:honest-bookkeeping}) and
the repaired alignments of Remark~\ref{rem:repaired-alignments}, for
every pair of primitive roles \((P_i, P_j)\) with \(i \ne j\), the
``forbidden collapse'' of the corresponding row in
Table~\ref{tab:boundary-matrix} is not admissible as an identification
of \(P_i\) with \(P_j\); the ``permitted interaction'' of the same row is
not thereby excluded and may remain admissible when its own bookkeeping
conditions are met.
\end{proposition}

\begin{remark}[Typed non-collapse versus global independence]\label{rem:typed-non-collapse-vs-independence}
Typed non-collapse is strictly weaker than global primitive
independence. It does not assert that the six roles are logically or
categorically independent in a global algebra; it asserts only that
the pairwise identifications listed as forbidden collapses in
Table~\ref{tab:boundary-matrix} are inadmissible under honest
bookkeeping, leaving the corresponding typed interactions available.
\end{remark}


\subsection{Level-profile architecture}\label{subsec:level-profile-architecture}

Typed non-collapse fixes which identifications are inadmissible; the
regime must also record the level at which each admissible claim is
made. The level trichotomy of Definition~\ref{def:level-trichotomy}
is therefore lifted from a role-by-role template to a regime-level
structure.

\begin{proposition}[Level-profile architecture]\label{prop:level-profile}
Every admissible claim about a description \(D \in \FATCD{}\) carries a
level profile: each claim is located on the behavioral, verification, or
structural-categorical level, and each cross-level assertion is
accompanied by a translation record in the sense of
Proposition~\ref{prop:honest-bookkeeping}. Direct comparison of claims at
different levels without a translation record is inadmissible.
\end{proposition}

\subsection{Forgetting and selected lifts}\label{subsec:forgetting-lifts}

Level transitions are controlled by forgetting and selected lifts.

\begin{proposition}[Forgetting and selected-lift architecture]\label{prop:forgetting-lifts}
Forgetting from the structural-categorical level to verification, and
from verification to behavioral, is admissible and is always lossy in
principle: an admissible forgetting carries a record of the losses in
\(\Audit{D}\). Lifting from a lower level to a higher level is
admissible only when accompanied by
\begin{itemize}[topsep=2pt,itemsep=1pt]
  \item the added data required to support the higher-level statement;
  \item the mark that the lift is selected and non-unique;
  \item explicit preservation of lower-level invariants and explicit
    losses where the lift is partial.
\end{itemize}
Lossless forgetting and unique lifts are inadmissible unless the losslessness
or uniqueness is itself proved and recorded.
\end{proposition}

\begin{remark}[Lower-to-higher non-determinacy]\label{rem:lower-to-higher-non-determinacy}
Proposition~\ref{prop:forgetting-lifts} formalizes the non-determinacy of
lower-level data: given behavioral or verification content alone, no
choice of structural-categorical content is forced by the admissibility
regime, and any asserted choice must be marked as a selected lift.
\end{remark}

\subsection{Activation thresholds}\label{subsec:activation-thresholds}

The regime must also state when a projected role is active. The
attachment protocol of Definition~\ref{def:threshold-attachment} is
lifted to a regime-level principle.

\begin{proposition}[Activation thresholds]\label{prop:activation-thresholds}
An admissible active derived role projection \(\pi_i(D)\) exists only when
\(D\) carries the threshold, witness, level, collapse, and audit data
required by Definition~\ref{def:threshold-attachment}. Below-threshold
clusters carry an explicit \(\belowthr\) status, and boundary-collapsed
clusters carry \(\collapsedbc\). Arbitrary enrichment, adding structure
solely to force a role projection, is inadmissible unless the additions
are themselves honestly bookkept as raw features or annotations.
\end{proposition}

\subsection{Instrument-relative visibility and suppression}\label{subsec:visibility-suppression}

Threshold-governed activation still leaves open what an instrument
makes visible or suppresses. Admissibility is therefore
instrument-relative.

\begin{proposition}[Instrument-relative visibility]\label{prop:visibility}
Every admissible claim carries an explicit visibility record comprising
its visible content, its suppressed content, the instruments under which
the visibility and suppression are asserted, and the resulting
instrument-relative claim strength. Overreading a suppressed claim as
visible under a different instrument is inadmissible. Partial visibility,
in which a claim is visible only up to some coarsening or translation, is
admitted when the visibility record itself makes the partiality explicit.
\end{proposition}

\begin{remark}[No instrument-free truth]\label{rem:no-instrument-free}
Proposition~\ref{prop:visibility} does not claim that instruments are
free of distortion or that any single instrument is canonical. It
asserts only that admissible claims carry an instrument-relative
visibility record. Instrument completeness is deferred to
Section~\ref{sec:future-work}.
\end{remark}

\subsection{Empirical-bridge admissibility}\label{subsec:empirical-bridge}

Any claim crossing from the closure-algebraic side of the framework
to a substrate-level observable must pass through the
empirical-bridge discipline, related to but narrower than
the empirical-adequacy view of theories~\citep{vanFraassen1980}.

\begin{proposition}[Empirical-bridge admissibility]\label{prop:empirical-bridge}
An empirical-bridge claim about \(D \in \FATCD{}\) is admissible only when
it carries, in \(\Audit{D}\),
\begin{itemize}[topsep=2pt,itemsep=1pt]
  \item the substrate and observable to which the bridge relates;
  \item the instrument under which the observable is measured;
  \item the level map between the closure-algebraic content and the
    substrate-level content;
  \item a threshold-witness relation recording when the bridge is
    crossed;
  \item the nonclaims suppressed by the bridge record, in particular any
    structural content of \(D\) that is not tracked by the substrate or
    instrument.
\end{itemize}
Empirical realization --- the claim that \(D\) is itself a substrate
observation --- is stronger than empirical-bridge admissibility and is
not asserted by any empirical-bridge claim.
\end{proposition}


\subsection{Nonclaim closure across the seven schemas}\label{subsec:nonclaim-closure}

The seven schemas of
Propositions~\ref{prop:honest-bookkeeping}--\ref{prop:empirical-bridge}
together preserve the global nonclaims of
Remark~\ref{rem:atlas-global-nonclaims} and block a corresponding
class of overclaims.

\begin{proposition}[Admissibility-regime nonclaim closure]\label{prop:nonclaim-closure}
Under the seven schemas, the following remain non-claims of the
admissibility regime:
\begin{itemize}[topsep=2pt,itemsep=1pt]
  \item no full six-primitives algebra is claimed;
  \item no global primitive-independence theorem is claimed;
  \item no empirical realization is claimed;
  \item no unique lift theorem is claimed;
  \item no lower-level-to-higher-level determinacy is claimed;
  \item no instrument-free truth claim is made;
  \item direct cross-level comparison remains inadmissible without
    translation;
  \item forgetting remains lossy unless losslessness is explicitly proved;
  \item empirical-bridge claims require threshold data, level maps,
    instrument visibility records, and suppressed-structure nonclaims.
\end{itemize}
Each schema blocks a specific class of overclaims. Honest bookkeeping blocks
unrecorded claims, promotions, translations, lifts, visibility claims, and
empirical-bridge assertions; typed non-collapse blocks forbidden primitive
collapses and flat independence overclaims while leaving restricted typed
boundary interactions available; level-profile architecture blocks direct
cross-level comparison without translation; forgetting-and-lift architecture
blocks lossless forgetting, unique lifts, and lower-to-higher determinacy;
activation thresholds block arbitrary enrichment and threshold-free
nontriviality claims; instrument-relative visibility blocks instrument-free
truth and visible-to-suppressed overreads; empirical-bridge admissibility
blocks empirical-realization overclaims and bridge claims without threshold or
level-map data.
\end{proposition}
